\section{What exact agreement across three interfaces establishes}
\label{exec16:semantic-parity}

The common operation service supplies a useful starting point for a banking
analyst, a command-line workflow and an assistant using MCP. Each should retain
the same question, assumptions, result and limitations. That objective needs a
precise boundary. Here the library entry is \texttt{Runtime.invoke}; a lower
level numerical helper is not automatically the same public operation. A CLI
exit code, an MCP protocol error and a canonical scientific disposition also
belong to different layers. Requiring identical bytes across those layers
would obscure their meanings rather than preserve them.

For a request $q$ that is admissible through all three negotiated surfaces,
let $L(q)$ denote the canonical library response, $C(q)$ the decoded CLI
response, and $M_s(q)$ and $M_t(q)$ the structured and decoded text responses
from MCP. The desired equality is
\[
 L(q)=C(q)=M_s(q)=M_t(q),\qquad q\in\mathcal Q_{\rm shared}.
\]
Equality includes context identity, operation, execution status, scientific
disposition, result kind, units, set topology, qualification, evidence and
error. A normalization that discards the error or reduces an unbounded set to
a convenient point would not meet this requirement. The shared domain is
itself negotiated: a generic local dispatcher can represent an unsupported
operation that the MCP server deliberately does not advertise.
\citep{exec16_parity_observed,exec16_parity_design}

\subsection{Actual local observations, with independent result meanings}

A deterministic local comparison exercised thirteen scientific requests
through the library, CLI and real Python MCP client/server over stdio. The
server negotiated protocol 2025-11-25 using SDK 1.27.0. All thirteen canonical
responses agreed exactly, including MCP's structured and serialized-text
copies. The cases comprised capabilities, graph syntax, a supplied-DAG
back-door expression, unsupported dependent sampling, a contradicted premise,
an unsupported estimation method, four deterministic sensitivity-set shapes,
weak-moment inversion, a changed-context failure and observation-clock
diagnostics. Three further lifecycle calls exercised report submission,
terminal status and cancellation after completion; together the selected
cases exercised all nine advertised operations. This is actual local execution
on generated inputs, not a mock transport or a statistical performance study.
\citep{exec16_parity_observed}

The equality check is useful because the expected meanings were also reviewed.
The four sensitivity examples solve $b^2\leq4$, $b^2+1\leq0$,
$-b^2-1\leq0$ and $-b^2+4\leq0$. Their sets are respectively
$[-2,2]$, the empty set, the whole line and
$(-\infty,-2]\cup[2,\infty)$. Their uncertainty is explicitly
nonprobabilistic. Matching copies of these objects cannot turn them into
sampling confidence regions. The moment example retains its conditional
Gaussian-moment interpretation. The graph example retains the supplied graph
and the limitation that syntax review does not prove causal truth. The clock
example flags both an observation available exactly at the decision time and
one without an availability time under the requested strict-before rule.
\citep{exec16_parity_observed}

For banking use, these distinctions matter when the same request is reused
in a batch report or presented by an assistant. The dependent-sampling and
unsupported-method examples returned a completed invocation with an
unsupported scientific disposition. The contradicted-premise example returned
a completed invocation with refusal. All three had CLI exit zero and
\texttt{isError=false} in MCP; none produced the requested causal estimate.
Changing the sign convention while retaining a moment artifact bound to the
old context instead produced a failed canonical response, CLI exit four and
\texttt{isError=true}. The preserved context hash and failure prevent a
response about one contrast from being relabeled as the opposite contrast.
This is agreement about the supported boundary, not a general guarantee that
every possible scientifically inconsistent context is recognized.
\citep{exec16_parity_observed}

The shared report job preserved computed, unsupported and refused input
declarations, while labeling its verification as schema-only. Its submission
response was a queued handle; the terminal state belonged inside the subsequent
job-status payload. Cancelling the completed job left its completed report
intact. A proposed consumer must display those states separately. A finished
report is not independent scientific approval of its contents, and a successful
cancel call cannot be presented as reversal of an already completed action.
This local case covers cancellation after completion; it does not replace
the distinct recovery, in-flight cancellation and target-platform evidence.
\citep{exec16_parity_observed}

\subsection{Capability negotiation and errors are part of the scientific contract}

The unimplemented discovery request exposes the boundary of the equality
statement. Library and CLI returned the same canonical unsupported response,
with CLI exit zero. Calling the unadvertised discovery tool through MCP
produced JSON-RPC error $-32602$, ``Unknown causal tool.'' This follows the
adapter's separation between advertised operations and unimplemented canonical
requests. The protocol error says that this tool call is unavailable; it does
not say that a causal effect is nonidentified. A proposed client should retain
the negotiated capability set and resolve this availability question before
dispatch, using a common user-facing explanation while preserving the actual
wire error if a call is nevertheless made. It should not invent a scientific
response that the service never returned.
\citep{exec16_parity_observed,v12_mcp_legacy_tools,exec16_parity_design}

The invalid-schema case also behaved differently at the transport boundary.
An extra request property raised a library validation exception, produced
CLI \texttt{invalid\_request} with exit two, and produced an MCP tool error
without a canonical structured response. This is appropriate input rejection
for that case. The corrupt-replay case shows why a matching coarse error label
is not enough. A valid capability request was first saved under an idempotency
key; its local replay JSON was then deliberately malformed. The library raised
\texttt{JSONDecodeError}. The CLI returned \texttt{invalid\_request} with
exit two, while MCP returned ``Invalid local causal request'' and no canonical
response. The request had not changed. The failing bytes belonged to stored
execution state, so attributing the failure to the analyst's request is wrong.
This extends the earlier CLI failure finding with a real MCP observation.
The original replay bytes were restored after the isolated local trial.
\citep{exec16_parity_observed}

The source explains the attribution. Replay parsing occurs inside the runtime,
but JSON decoding errors inherit from \texttt{ValueError}; the CLI and MCP
adapters catch that broad exception family together with input-validation
errors. The proposed repair is to separate validated caller-input failures
from artifact, replay-state, worker and transport failures. For an already
validated request with a corrupt replay, return a failed/unavailable canonical
result with its context identity and a specific stored-state error. For an
unparseable incoming message, retain a transport/input error without fabricating
a scientific context. Preserve the corrupt record for operator investigation;
do not automatically discard an idempotency record and repeat a potentially
consequential action. This repair is specified, not implemented in the observed
runtime. \citep{exec16_parity_observed,exec16_parity_design}

\subsection{Validation ownership, identity and the remaining acceptance cases}

Validation responsibility cannot be assigned merely by noticing that a tool
advertises an output schema. In the inspected SDK, a handler returning an
explicit \texttt{CallToolResult} takes a path that returns that object before
the decorator's later structured-output validation branch. The local adapter
does explicitly validate its canonical response, and the runtime also checks
response semantics. A future refactor must preserve those checks and test
malformed responses on every return path. The successful local trials do not
establish that the SDK will perform validation which this specific code path
leaves to the adapter. \citep{exec16_parity_observed}

Execution identity needs an equally precise scope. The runtime binds selected
local sources, two schemas, Python and four numerical/validation dependencies.
Its identity does not itself include MCP or AnyIO versions or the separate
MCP tool-mapping contract. The current observations record SDK versions and
source identities separately. The proposed transport identity should bind
the negotiated protocol, tool map, SDK, concurrency dependency, launch profile
and any response normalizer as well as the scientific runtime. A common
scientific result with a changed transport contract then remains distinguishable
from an exact replay under the same complete interface.
\citep{exec16_parity_observed,exec16_parity_design}

The known next tests are concrete: valid requests with missing, malformed or
incompatible replay records; unknown tools and unsupported method identifiers;
schema and context mismatches; nonfinite or incorrectly typed responses;
divergent MCP text and structured content; changed SDK/tool-map identity;
and terminal versus in-flight job cancellation. Freeze the request and expected
layer of failure before each test, retain failed attempts, and compare the
whole result object. Extend the common decoder to each proposed consumer
language and then test against its actual caller. These engineering routes
are known, but their completion and external acceptance are not supplied by
the thirteen matched examples. \citep{exec16_parity_design}

The environment observation is also bounded. A six-second worker-thread probe
stalled under the managed sandbox; the same fixed local probe succeeded in
the approved execution environment outside it, where the MCP comparison ran.
This is an observed operating constraint, not a reason to treat sandbox and
production environments as interchangeable. HTTP, multiple operators,
authentication, remote retries, newer protocol revisions and bank-platform
performance remain separate qualification work. The immediate proposal can
justify testing these interfaces with generated and permitted public inputs;
actual consumer owners must later accept their deployed workflows. No transport
comparison can supply missing intervention truth, causal assumptions or useful
bank precision. \citep{exec16_parity_observed,exec16_parity_design}
\label{exec16:semantic-parity-end}
